\section{Discussion}
\label{sec:discussion}

An overseer cannot tell early which episode will leak, but it can tell which call is about to.
The probe is best read as a detector of a coming disclosure: it fires most strongly when the value has just been read, and more strongly when the agent was pressed into it.
It does not read the rule text, so a disclosure right after a read is flagged even where it is allowed (\cref{app:e15}).
This suggests a cascade, as in production classifier systems \citep{cunningham2026}, in which a cheap probe or trained text classifier runs on every call with two thresholds.
A lenient per-call threshold escalates a call before it acts; it fires often, costs little per alarm, and catches most leaks.
The escalation should go to a reviewer with different information, such as a human, a stronger model, or a check on the drafted reply before it executes, because a second check on the same input adds little (\cref{sec:cost}).
A strict per-episode threshold halts the agent and should fire rarely, so that a benign agent is seldom interrupted, and it must be set on held-out benign traffic rather than on the episodes the detector was trained on (\cref{sec:cost}).
Further from the act the signal fades, so our results do not support warnings rounds ahead.
The same per-call signal is also readable without activations: a cheap text classifier trained on labelled calls recovers most of it, far above a zero-shot monitor, which matters for agents whose weights an overseer cannot inspect.
Whether that holds outside templated scenarios like ours is open.

\subsection*{Limitations}
We list the main limitations in order of how much each could change our conclusions; \cref{app:limitations} gives the rest.

\paragraph{Narrow setting.}
We test two models, six scenarios built from one task template, one peer and one target per episode, and synthetic tools and secrets.
Leaks are exact strings, which an output filter would catch. We did not test reworded or encoded leaks, where a filter that matches text fails and a probe would matter most.
Qwen3-32B ran with its reasoning switched off.
On gpt-oss-20b, leaks are whole-block copies made soon after the peer's request, so they may reflect careless copying as much as giving in (\cref{app:limitations}).
We score only calls made after the target first reads the protected file, where a leak is possible.
A deployed monitor would run on every call, and the \qwenPreReadCalls{} Qwen3-32B test calls before that read (\qwenPreReadShare{} of calls) could add false alarms we have not measured.
Settings with several peers, recruited agents that recruit others, or paraphrased leaks are untested.

\paragraph{Few leaking test calls on gpt-oss-20b.}
On gpt-oss-20b, the committed test covers one probe, and its test episodes hold few leaking calls (\gptossFreshLeaks{}), so its interval is wide (\cref{app:confirmation}).

\paragraph{What the probe reads is open.}
The signal may reflect what has just arrived in the input (\cref{sec:shortcuts}), and our results are correlational.
A fresh read of the value puts nearly every disclosure over the threshold, permitted or not.
Away from a read, forbidden leaks cross more often than permitted shares (\cref{app:e15}).
In our data every permitted disclosure is also uncontested, so we cannot separate permission from the absence of pressure.

\paragraph{An untuned text monitor.}
Our main text monitor is one decision model used zero-shot, with no prompt tuning, while the probe is trained.

\paragraph{Future work.}
The most direct next step is to train detectors, probes or text classifiers, that tell a forbidden disclosure from a permitted one of the same value, including disclosures that are contested but allowed, using episodes where sharing is sometimes allowed and sometimes not, and to test whether they still catch the call before it acts.
Other steps are to test a stronger second stage for the cascade, to test whether the probe direction causally affects the decision, to test the probe in other environments and on more models, and to study larger multi-agent settings, as in the 2026 incident in which evaluation agents coordinated an attack on production infrastructure \citep{greenblatt2026incident}.
